In regular waves of height $\gls{sym-H}$, the conditions for slamming can be derived by comparing the vertical position of the bottom of the WEC, $\gls{sym-xi-t}-\gls{sym-Delta-z-slam}$, with the wave elevation $\gls{sym-zeta-y-t}$, which depends on the horizontal position coordinate $\gls{sym-y}$.
To prevent slamming, the criteria
\begin{equation}\label{eq:slamming-time-domain}
\begin{aligned}
    \gls{sym-xi-t} - \gls{sym-Delta-z-slam} &< \gls{sym-zeta-y-t} \\
    |\gls{sym-hat-xi}|\cos(\gls{sym-omega} \gls{sym-t}+\angle \gls{sym-hat-xi}) - \gls{sym-Delta-z-slam} &< \frac{\gls{sym-H}}{2}\cos(\gls{sym-omega} \gls{sym-t} - \gls{sym-k-wavenumber}\gls{sym-y})
\end{aligned}
\end{equation}
must be true at all times and over all positions where the body could exit the waves, $-\gls{sym-D}/2<\gls{sym-y}<\gls{sym-D}/2$.
Note that this expression assumes a sinusoidal waveshape but still applies for irregular waves if a wave-by-wave approach \citep{lin_fast_2025} is utilized.
Manipulating \Cref{eq:slamming-time-domain} to ensure the inequality is satisfied for all $\gls{sym-t}$ yields the criterion:
\begin{equation}
|\gls{sym-hat-xi}|^2 - \gls{sym-H}\cos(\gls{sym-k-wavenumber}\gls{sym-y})|\gls{sym-hat-xi}|\cos\angle\gls{sym-hat-xi} + \gls{sym-H}\sin(\gls{sym-k-wavenumber}\gls{sym-y})|\gls{sym-hat-xi}|\sin\angle\gls{sym-hat-xi} < \gls{sym-Delta-z-slam}^2-\left(\frac{\gls{sym-H}}{2}\right)^2
\end{equation}
which is quadratic in $\gls{sym-hat-xi}$ and therefore already in a form suitable for the constrained optimal control procedure described in \Cref{sec:optimal-control,sec:appendix-qp-solution}, once the appropriate worst-case $\gls{sym-y}$ is substituted.
The remainder of this section develops an explicit expression for the allowable amplitudes under various scenarios and the corresponding worst-case $\gls{sym-y}$.
Completing the square results in:
\begin{equation}\label{eq:slamming-squared}
    \left[ |\gls{sym-hat-xi}|- \frac{\gls{sym-H}}{2}\cos(\angle\gls{sym-hat-xi} + \gls{sym-k-wavenumber}\gls{sym-y})\right]^2 < \gls{sym-Delta-z-slam}^2 - (\gls{sym-H}/2)^2 \sin^2(\angle\gls{sym-hat-xi} + \gls{sym-k-wavenumber}\gls{sym-y})
\end{equation}
In general, after taking the square root of \Cref{eq:slamming-squared}, both the positive and negative branches must be used.
This yields a maximum and minimum amplitude criteria to prevent slamming, $\gls{sym-xi}_{\text{min},\text{slam}}<|\gls{sym-hat-xi}|<\gls{sym-xi}_{\text{max},\text{slam}}$, with the following values:
\begin{equation}\label{eq:slamming-min-max}
\begin{aligned}
    \xi_{max,slam} &= \frac{H}{2} \cos\theta + \sqrt{\Delta z_{slam}^2-\left(\frac{H}{2}\sin\theta\right)^2 } \\
    \xi_{min,slam} &= \frac{H}{2} \cos\theta - \sqrt{\Delta z_{slam}^2-\left(\frac{H}{2}\sin\theta\right)^2 }
\end{aligned}
\end{equation}
where $\gls{sym-theta}=\angle\gls{sym-hat-xi} + \gls{sym-k-wavenumber}\gls{sym-y}$ must be evaluated with the value of $\gls{sym-y}$ that creates the minimum value of $\gls{sym-xi}_{\text{max},\text{slam}}$ and the maximum value of $\gls{sym-xi}_{\text{min},\text{slam}}$ to ensure lack of slamming across the entire WEC surface.
The relevant value of $\gls{sym-theta}$ differs depending on the wave amplitude.

%%%%%%%%%%%%%%%
\subsubsection{Case 1: Small Wave Amplitudes}
For the case of $\gls{sym-Delta-z-slam} >\gls{sym-H}/2$ (wave amplitudes smaller than the draft), \Cref{eq:slamming-min-max} reduces to only a maximum amplitude criteria $|\gls{sym-hat-xi}|<\gls{sym-xi}_{\text{max},\text{slam}}$ because $\gls{sym-xi}_{\text{min},\text{slam}}<0$.
This aligns with the intuition that a body fixed at the still water line would not exit the water when the wave amplitude is less than the draft.
The value of $\gls{sym-theta}$ that minimizes $\gls{sym-xi}_{\text{max},\text{slam}}$ over the relevant range is
\begin{equation}\label{eq:slamming-limit-small-waves}
\theta = \pi - \max\left(0,~ \frac{-\gls{sym-k-wavenumber}\gls{sym-D}}{2}+|\pi-\angle \xi|\right)
\end{equation}
where the symmetry of \Cref{eq:slamming-min-max} has been used to collapse possible values onto $[0,\gls{sym-pi}]$.

This expression reduces to the simple limits $\gls{sym-Delta-z-slam} -\gls{sym-H}/2$ and $\gls{sym-Delta-z-slam} +\gls{sym-H}/2$ for short wavelengths $(\gls{sym-theta}=\gls{sym-pi})$ and long wavelengths $(\gls{sym-theta}=0)$ respectively.
These limits are also the minimum and maximum values of $\gls{sym-xi}_{\text{max},\text{slam}}$ for any diameter and wave condition where $\gls{sym-Delta-z-slam} >\gls{sym-H}/2$.
For intermediate wavelengths, the variable $\gls{sym-theta}$ accurately accounts for the effect of body diameter and the phase offset of body motion from the waves.
Intuitively, $\gls{sym-theta}$ represents how high up the free surface the bottom of the WEC is allowed to get before the edge exits the water, with $\gls{sym-theta}=0$ allowing the body to get all the way to the wave crest (top) and $\gls{sym-theta}=\gls{sym-pi}$ requiring the float to remain fully below the wave trough (bottom).

%%%%%%%%%%%%%%%
\subsubsection{Case 2: Large Wave Amplitudes}
In sufficiently large waves ($\gls{sym-H}/2>\gls{sym-Delta-z-slam}$), \Cref{eq:slamming-limit-small-waves} does not apply.
Intuitively, when the wave amplitude exceeds the draft, a minimum amplitude of motion in-phase with the incident wave is required to prevent the wave trough from going below the body surface.
The minimum amplitude arises mathematically when $\gls{sym-xi}_{\text{min},\text{slam}}>0$, which occurs exactly for waves meeting the amplitude criterion.
The ``in-phase'' motion requirement is consistent with the observation that in this wave regime, the maximum amplitude in \Cref{eq:slamming-min-max} becomes unsatisfiable ($\gls{sym-xi}_{\text{max},\text{slam}}<0$) for WEC phases that are more than 90 degrees out of phase with the incident wave ($\gls{sym-theta} \\text{in} [\gls{sym-pi}/2,3\gls{sym-pi}/2]$).

For evaluating \Cref{eq:slamming-min-max}, the value of $\gls{sym-theta}$ that produces both minimum $\gls{sym-xi}_{\text{max},\text{slam}}$ and maximum $\gls{sym-xi}_{\text{min},\text{slam}}$ in this regime is:
\begin{equation}
\theta = \angle\hat{\xi} + \frac{\gls{sym-k-wavenumber}\gls{sym-D}}{2}\textrm{sgn}(\pi-\angle\hat{\xi})
\end{equation}
After wrapping the value onto $[0,2\gls{sym-pi}]$, it can be mapped to $[0,\gls{sym-pi}]$ without changing the value of \Cref{eq:slamming-min-max} via the identity $\cos(\gls{sym-theta})=\cos(\gls{sym-pi}-|\gls{sym-pi}-\gls{sym-theta}|)$.

Notably, in this wave regime it becomes possible for the right hand side of \Cref{eq:slamming-squared} to become negative for certain values of $\angle\gls{sym-hat-xi}$ and $\gls{sym-D}$.
This means that those combinations of $\angle\gls{sym-hat-xi}$ and $\gls{sym-D}$ are prohibited because slamming would occur regardless of the amplitude $|\gls{sym-hat-xi}|$.
The following requirement on $\gls{sym-D}$ ensures that there exist some values of $\angle\gls{sym-hat-xi}$ without automatic slamming:
\begin{equation}\label{eq:D-slamming}
    \gls{sym-k-wavenumber}\gls{sym-D}<2\arcsin\left(\frac{\Delta z_{slam}}{H/2}\right)
\end{equation}
which has upper bound $\gls{sym-k-wavenumber}\gls{sym-D}<\gls{sym-pi}$.
Once \Cref{eq:D-slamming} is satisfied, the set of acceptable $\angle\gls{sym-hat-xi}$ that avoid negative values for both $\gls{sym-xi}_{\text{max},\text{slam}}$ and the right hand side of \Cref{eq:slamming-squared} are:
\begin{equation}\label{eq:slamming-allowed-angles}
\begin{aligned}
    \angle\hat{\xi} \in &\left[0,~~\arcsin\left(\frac{\Delta z_{slam}}{H/2}\right)-\frac{\gls{sym-k-wavenumber}\gls{sym-D}}{2}\right]
    \cup \left[2\pi-\arcsin\left(\frac{\Delta z_{slam}}{H/2}\right)+\frac{\gls{sym-k-wavenumber}\gls{sym-D}}{2}, ~~2\pi\right]
\end{aligned}
\end{equation}
At the edge-case $\angle\gls{sym-hat-xi}$ on the inner boundaries of this interval, only a single amplitude is permitted because the values for $\gls{sym-xi}_{\text{max},\text{slam}}$ and $\gls{sym-xi}_{\text{min},\text{slam}}$ coincide at $\sqrt{(\gls{sym-H}/2)^2-\gls{sym-Delta-z-slam} ^2}$.
At any lower amplitude, the WEC would exit the water at the wave trough, and any higher, the WEC would raise out of the water at the wave crest.
The permissible amplitude range can be widened by decreasing $\gls{sym-D}$ or moving $\angle\gls{sym-hat-xi}$ towards the closer of $0$ and $2\gls{sym-pi}$.
The widest permissible amplitude range occurs when the WEC is perfectly in-phase with the incident wave ($\angle\gls{sym-hat-xi}=0$) and is small with respect to the wavelength ($\gls{sym-k-wavenumber}\gls{sym-D}\approx0$), which yields the limits $\gls{sym-H}/2-\gls{sym-Delta-z-slam} < |\gls{sym-hat-xi}| < \gls{sym-H}/2+\gls{sym-Delta-z-slam}$.

%%%%%%%%%%%%%%%
\subsubsection{Implementation}
\Cref{fig:slam} visualizes the nondimensional maximum and minimum slamming amplitudes $\gls{sym-xi}_{\text{max},\text{slam}}/(\gls{sym-H}/2)$ and $\gls{sym-xi}_{\text{min},\text{slam}}/(\gls{sym-H}/2)$ as a function of the worst-case phase angle $\gls{sym-theta}$ and the draft to wave amplitude ratio $\gls{sym-Delta-z-slam} /(\gls{sym-H}/2)$.
This reveals that besides increasing the draft, decreasing $\gls{sym-theta}$ can prevent slamming, for example by decreasing the diameter or adjusting $\angle \gls{sym-xi}$ via control.
If $\gls{sym-Delta-z-slam} <\gls{sym-H}/2|\sin\gls{sym-theta}|$, the design is unviable and the slamming limits are undefined, consistent with the negative right hand side of \Cref{eq:slamming-squared} discussed above.
\begin{figure}[htbp]
    \centering
    \includegraphics[\FigureOptions{aor:figs/from-matlab/slamming_amplitude.pdf}]{figs/from-matlab/slamming_amplitude.pdf}
    \caption{Nondimensional critical slamming amplitude}
    \label{fig:slam}
\end{figure}

Separately, there exists a symmetrical counterpart to the slamming condition that we will call the submersion condition.
While preventing slamming requires that the bottom of the WEC remains below the free surface, preventing submersion requires that the top of the WEC remains above the free surface.
Analogous to \eqref{eq:slamming-time-domain}, we obtain:
\begin{equation}\label{eq:submersion}
    \gls{sym-xi-t} + \gls{sym-Delta-z-sub} > \gls{sym-zeta-y-t}
\end{equation}

Conceptually, both the slamming requirement \eqref{eq:slamming-time-domain} and the submersion requirement \eqref{eq:submersion} must be applied to both the float and the spar for a total of 4 free surface constraints for each sea state.
However, the limit on $|\gls{sym-hat-xi-f}-\gls{sym-hat-xi-s}|$ described in \Cref{tab:DLCs} already ensures that the top of the spar remains above the top of the float across all dynamic conditions.
Coupled with the fact that the float diameter always exceeds the spar diameter, this means that satisfaction of the float submersion requirement implies automatic satisfaction of the spar submersion requirement.
While this logic works for submersion, it does not apply for slamming because the damping plate diameter may exceed the float diameter.
The remaining three free surface constraints are aggregated into a single constraint for each sea state, using the minimum z-dimension for each body $\gls{sym-Delta} \gls{sym-z}_{\text{surf}}=\\text{min}(\gls{sym-Delta-z-sub}, \gls{sym-Delta-z-slam} )$, where the values for $\gls{sym-Delta-z-sub} $ and $\gls{sym-Delta-z-slam}$ were given in \Cref{sec:dynamics}.
This reduces the number of free surface constraints per sea state from four to one.

For operational waves, the vertical dimension $\gls{sym-Delta} \gls{sym-z}$ in \Cref{eq:slam-constraint-quadratic} is set as follows.
For float slamming, $\gls{sym-Delta} \gls{sym-z} = \gls{sym-T-f-1}$ (rather than $\gls{sym-T-f-2}$): this prevents slamming on the slanted underside of the float and additionally maintains constant waterplane area, avoiding unmodeled nonlinearities in the hydrostatic stiffness and Froude–Krylov excitation force.
For spar slamming, $\gls{sym-Delta} \gls{sym-z} = \gls{sym-T-s} - \gls{sym-h-d}$ prevents both slamming on the bottom of the spar and surfacing of the damping plate.
For float and spar submersion, $\gls{sym-Delta} \gls{sym-z}$ is set to the above-water heights $\gls{sym-h-f} - \gls{sym-T-f-2}$ and $\gls{sym-h-s} - \gls{sym-T-s}$ respectively.

The storm case slamming constraint is not currently applied.
A future extension should set $\gls{sym-Delta} \gls{sym-z} = \gls{sym-T-f-2}$ in the storm case to prevent slamming on the (now-merged) bottom surface of the float-spar system.

The amplitude limits $\gls{sym-h}_{fs,\text{clear}}$, $\gls{sym-h}_{fs,\text{up}}$, $\gls{sym-h}_{fs,\text{down}}$, $\gls{sym-xi}_{\gls{sym-f},\text{linear}}$, and $\gls{sym-xi}_{\gls{sym-s},\text{linear}}$ are sea-state-independent upper limits, while $\gls{sym-xi}_{\text{slam}}$ depends on the sea state and can act as both an upper and lower limit on oscillation magnitude.

Relevant model values for the nominal float design under reactive optimal control are plotted for each sea state in \Cref{fig:slamming-validation}.
Comparing the first two subplots for the phase of the float motion $\angle\gls{sym-hat-xi-f}$ and the slamming angle $\gls{sym-theta}$ shows that the slamming angle roughly ``wraps'' the motion angle around a value of $\gls{sym-pi}/2$, colored white in the figure.
The third plot shows that the minimum slamming and submersion amplitude is not a concern since it is negative for all sea states and thus will never be an active constraint.
The fourth plot shows the maximum amplitude limit and the fifth depicts the actual amplitude, with amplitudes in violation of the limit indicated with an overlaid black hatching.
A significant fraction of sea states violate the limit and will require adjusted control gains computed via constrained optimal control per \Cref{sec:optimal-control,sec:appendix-qp-solution}.
Significant wave heights in excess of 5.5 meters have negative maximum slamming amplitudes (colored blue), indicating that control adjustments in those waves must not only decrease $|\gls{sym-hat-xi-f}|$ but also adjust $\angle\gls{sym-hat-xi-f}$ to reduce $\gls{sym-theta}$.
\begin{figure}[htbp]
    \centering
    \includegraphics[\FigureOptions{aor:figs/from-matlab/slamming_model_comparison_float.pdf}]{figs/from-matlab/slamming_model_comparison_float.pdf}
    \caption{Slamming Model Results for Nominal RM3 Float Design in Operational Sea States}
    \label{fig:slamming-validation}
\end{figure}